\subsection{Unramified ring maps and the exact diagonal idempotent}\label{algebra-proof-056-unramified-ring-maps-and-the-exact-diagonal-idempotent}

\subsubsection{Source and scope}\label{algebra-proof-056-source-and-scope}

Source: Stacks Project authors, algebra.tex at commit a044\allowbreak{}46e5\allowbreak{}7ec1\allowbreak{}fbc2\allowbreak{}52a8\allowbreak{}71af\allowbreak{}cec7\allowbreak{}752f\allowbreak{}b280\allowbreak{}7b14, SHA256 FA8B\allowbreak{}B92E\allowbreak{}58A4\allowbreak{}F78A\allowbreak{}2BD0\allowbreak{}1B3B\allowbreak{}6A4A\allowbreak{}87DE\allowbreak{}0A0D\allowbreak{}279F\allowbreak{}5DD9\allowbreak{}0641\allowbreak{}B574\allowbreak{}DD5F\allowbreak{}BFFF\allowbreak{}A4F3. The full 41377--41711 interval and twelve received reports are read here; 41712--41723 is read but unadjudicated. Dependency reading is bounded to 539--567, 2427--2458, 3921--3953, 38608--38652, 32653--32686, 27420--27450 and 39393--39419. These bounded excerpts are not claims that all the associated proofs have been newly read. The already reviewed formal, smooth, étale and local-integral notes retain their prior evidence. Four indexed hits for ``unramified diagonal'' remain unread and unused. The source's Henselian citation and EGA terminology are retained; no new reading of those publications or novelty claim is made.

The source calls finite type plus \(\Omega_{S/R}=0\) unramified, and finite presentation plus the same condition G-unramified. Keep these definitions exactly. Neither a translation nor this review silently changes them to a later preferred convention. The fuller arguments and diagonal consequences below are separate editorial material.

\subsubsection{Source statements and the two reference repairs}\label{algebra-proof-056-source-statements-and-the-two-reference-repairs}

Formal unramifiedness is equivalent to \(\Omega_{S/R}=0\), as proved in the preceding section, so 41404--41421 follows with exactly the two different finiteness hypotheses. For the thirteen assertions at 41423--41547, base change uses the differential base-change isomorphism and the stated base-change finiteness lemma. Composition instead uses \[
T\otimes_S\Omega_{S/R}\longrightarrow\Omega_{T/R}
\longrightarrow\Omega_{T/S}\longrightarrow0
\] and lemma-\AlgebraIdentifierBreak{}compose-\AlgebraIdentifierBreak{}finite-\AlgebraIdentifierBreak{}type, which proves both composition finiteness statements at 539--563. The source reference at 41477 to the base-change finiteness lemma is replaced by this actual composition lemma; the maps and claim remain the same.

For a principal localization, every relative derivation kills \(R\) and kills \(f^{-1}\), since \(d(f^{-1})=-f^{-2}df=0\); the presentation \(R[X]/(fX-1)\) is finite. For \(R/H\) every relative derivation is zero, because it kills the surjective image of \(R\). The quotient is finite type and is finitely presented when \(H\) is finitely generated. Étale maps supply finite presentation and zero differentials as already proved.

For finite-type \(R\to S\), write \(S\) using finitely many generators \(x_1,\ldots,x_n\). Their differentials generate \(\Omega_{S/R}\). If its localization at \(\mathfrak q\) vanishes, choose one denominator outside \(\mathfrak q\) annihilating each \(dx_i\), and multiply those finitely many denominators to get \(g\notin\mathfrak q\) annihilating the module after localization. Then \(\Omega_{S_g/R}=0\). Finite type, and finite presentation when assumed, survive this principal localization. The residue-field criterion is equivalent by Nakayama for this finite localized module. Differential base change identifies the fibre criterion with the same residue vector space, so all four local tests in items (6)--(9) agree. Their finite-type hypothesis is used to pass from residue vanishing to an actual open neighbourhood.

For a finite principal cover \(D(g_j)\), vanishing of the localized differentials implies global vanishing; local finite type and local finite presentation glue by lemma-cover. This proves item (10). In item (11), neither global finiteness hypothesis has yet been given. Its hypothesis instead gives, at every prime, a principal open where the corresponding finite-type or finite-presentation condition and vanishing differentials hold. Quasi-compactness extracts finitely many such opens. They cover the spectrum, so the \(g_j\) generate the unit ideal, and item (10) applies. The reference to (6) and (7) at 41519 must therefore be replaced by the definition and (10). This is the actual argument, not an added global finiteness assumption. For \(S=0\), the map is already finitely presented as \(R/(1)\) and has zero differentials, so the empty-spectrum case is also covered.

For item (12), retain the exact finite presentation \(S=R[x_1,\ldots,x_n]/(g_1,\ldots,g_m)\). The vanishing of the differential cokernel gives all the actual identities \[
dx_i=\sum_j h_{ij}\,dg_j+\sum_{j,k}a_{ijk}g_j\,dx_k
\] in the original free differential module over \(R[x_1,\ldots,x_n]\). Adjoin to \(\mathbf Z\) every coefficient of \(g_j,h_{ij},a_{ijk}\) to get the subring \(R_0\subset R\). These finitely many identities hold over \(R_0\) itself, because its inclusion in \(R\) is injective. Thus the displayed \(S_0\) is finitely presented with \(\Omega_{S_0/R_0}=0\), and its original base change is \(S\). For item (13), use the source's full identities \[
dx_i=\sum_j h_{ij}\,dg_{ij}+\sum_k g'_{ik}\,dx_k,\qquad
g_{ij},g'_{ik}\in I.
\] Only finitely many of these polynomials occur. Keeping every coefficient gives the source \(S_0=R_0[x_1,\ldots,x_n]/(g_{ij},g'_{ik})\). It has zero relative differentials and its base change surjects onto the original \(S=R[x_1,\ldots,x_n]/I\). In fact the constructed \(S_0/R_0\) is already G-unramified, since its displayed relation list is finite. This is an explicit consequence of the existing proof, not an error in the source's weaker conclusion.

For 41567--41641, replace \(S\) only by the stated neighbourhood and then its actual fibre \(F=S\otimes_R\kappa(\mathfrak p)\). It is finite type over the field and has zero differentials. The field smoothness criterion at 38609 applies because zero is at most its nonnegative local component dimension; hence it is smooth with zero differentials, and is étale. The already proved field characterization identifies \(F\) with a finite product of finite separable fields. Its local ring at the original fibre prime is \[
F_{\bar{\mathfrak q}}=S_{\mathfrak q}/\mathfrak pS_{\mathfrak q}
=\kappa(\mathfrak q).
\] This gives both equality of maximal ideals and the exact finite separable residue extension. The fibre point is isolated, giving quasi-finiteness. Conversely, the two residue conditions give that same fibre local ring as a finite separable field; its differential module is zero, and the finite-type fibre test proves unramifiedness. No claim of flatness follows from these residue conditions alone.

The equivalence étale \(\Leftrightarrow\) flat plus G-unramified retains finite presentation. The previous étale criterion applies to the flat local maps with the just proved zero-dimensional separable fibres. The reverse implication uses the already proved flatness of étale maps.

Finally let the source map \(\varphi:k[x_1,\ldots,x_n]\to A,\ x_i\mapsto a_i\), be finite type. This is finite presentation because its polynomial base is Noetherian. If \(da_1,\ldots,da_n\) generate \(\Omega_{A/k}\), the exact differential sequence gives \(\Omega_{A/k[x_1,\ldots,x_n]}=0\), so \(\varphi\) is unramified. At a maximal ideal \(\mathfrak m\) of \(A\), its residue field is finite algebraic over \(k\). The image of the polynomial base in it is generated by finitely many algebraic elements; it is a finite dimensional domain over \(k\), hence a field. Thus \(\varphi^{-1}(\mathfrak m)\) is maximal, and its regular local polynomial ring has dimension \(n\). Under condition (a), the target local dimension is also \(n\); condition (b) bounds the differential residue dimension by \(n\). The field smoothness criterion makes \(A_{\mathfrak m}\) regular, hence Cohen--Macaulay. Its fibre over the polynomial local base is the finite separable residue field just proved, of dimension zero. All hypotheses of miracle flatness at 32654 are now met, including the exact dimension equality \(n=n+0\). This proves flatness at every maximal ideal and hence globally. Conversely, étaleness gives condition (b), flatness and the same zero-dimensional fibres. The original local dimension formula gives condition (a). For \(A=0\), the criteria are vacuous and the empty étale algebra still satisfies the statement; no maximal-ideal argument is needed.

\subsubsection{The diagonal, its determinant and equality of maps}\label{algebra-proof-056-the-diagonal-its-determinant-and-equality-of-maps}

For any map \(R\to S\), put \[
D=S\otimes_RS,\qquad \mu:D\to S,\quad a\otimes b\mapsto ab,\qquad
J=\ker\mu.
\] The exact diagonal-differential comparison gives \(J/J^2=\Omega_{S/R}\). When \(R\to S\) is finite type, \(J\) is generated by the actual elements \(u_i=x_i\otimes1-1\otimes x_i\), where the \(x_i\) are the chosen algebra generators. Indeed the quotient by these elements identifies the two structural maps from \(S\); the map \(S\to D/(u_i)\), \(s\mapsto s\otimes1\), is then inverse to \(\mu\). Thus an unramified map has a finitely generated ideal \(J=J^2\).

More generally only finite generation of \(J\), not finite type of \(S\), is needed in the following exact construction. Choose actual generators \(u_1,\ldots,u_m\) of \(J=J^2\), and write \[
u_i=\sum_{j=1}^m a_{ij}u_j,\quad a_{ij}\in J,\qquad
e=\det(I_m-(a_{ij}))
=\sum_{\pi\in\mathfrak S_m}\operatorname{sgn}(\pi)
\prod_{i=1}^m(\delta_{i,\pi(i)}-a_{i,\pi(i)}).
\tag{UR-1}
\] The adjugate applied to \((I_m-(a_{ij}))(u_i)_i=0\) gives \(eJ=0\). Every entry of the coefficient matrix lies in \(J\), so \(e-1\in J\), and \(e(e-1)=0\). Hence \[
e^2=e,\quad \mu(e)=1,\quad J=(1-e)D.
\tag{UR-2}
\] For the last equality, \(1-e\in J\) and each \(j\in J\) is \((1-e)j\). For \(m=0\), the empty determinant is \(1\), \(J=0\), and the same formulas hold. These formulas also hold for the zero ring.

The maps \[
S\longrightarrow eD,\quad s\longmapsto(s\otimes1)e,\qquad
eD\longrightarrow S,\quad z\longmapsto\mu(z)
\tag{UR-3}
\] are inverse: for \(d\in D\), \(d-\mu(d)\otimes1\in J\), annihilated by \(e\). The first map is \(D\)-linear when \(S\) is a \(D\)-module through \(\mu\), since that same difference is killed by \(e\). Consequently \(S\) is a direct summand of the rank-one free \(D\)-module \(D=eD\oplus(1-e)D\), hence is a finite projective \(D\)-module, and \(D_e\to S\) is the exact original multiplication quotient. Its inverse sends \(s\) to \((s\otimes1)/1\). In \(D_e\) the idempotent \(e\) is \(1\), so this agrees with (UR-3). This is a statement about the diagonal algebra \(D\), not projectivity or finiteness of \(S\) as an \(R\)-module.

The following three conditions on an arbitrary \(R\to S\) are equivalent: formal unramifiedness with \(J\) finitely generated; existence of \(e\) as in (UR-2); and projectivity of \(S\) as a \(D\)-module. The first implication is (UR-1), and the second gives projectivity by (UR-3). For the converse, a \(D\)-linear section \(\sigma\) of \(\mu\) exists. Put \(e=\sigma(1)\). Then \(\mu(e)=1\) and \(Je=0\). The same argument gives \(e^2=e\) and \(J=(1-e)D\), hence \(J\) is finitely generated and \(J/J^2=0\). Thus \(\Omega_{S/R}=0\). The idempotent is unique: if \(e'\) has the same properties, then \(e-e'\in J\), so \(e(e-e')=e'(e-e')=0\), giving \(e=ee'=e'\). The determinant construction is therefore independent of the selected generators and expressions, even though it preserves them as its explicit formula.

Every base change \(R\to R_1\) carries this idempotent to the diagonal idempotent for \(S_1=R_1\otimes_RS\), under \(R_1\otimes_RD\cong S_1\otimes_{R_1}S_1\). The multiplication map has the \(R\)-linear section \(s\mapsto s\otimes1\), so its kernel sequence remains exact after this base change. Alternatively the decomposition in (UR-2) itself tensors to the asserted kernel decomposition. Neither argument assumes \(R_1\) flat.

For two \(R\)-algebra maps \(f,g:S\to B\), let \(h:D\to B,\ h(a\otimes b)=f(a)g(b)\), and put \(d=h(e)\). The ideal generated by all differences \(f(s)-g(s)\) is exactly \[
\mathfrak a=h(J)B=(1-d)B,\qquad d^2=d.
\tag{UR-4}
\] Here \(J\) is generated by the differences in the two tensor factors, and (UR-2) gives the equality. Thus \(B/\mathfrak a\cong B_d\), and on this factor the two original maps are equal. It is the largest factor on which they agree: for every \(B\)-algebra \(B'\), equality of their composites is equivalent to the image of \(1-d\) being zero, equivalently the image of \(d\) being \(1\), equivalently factorization through \(B_d\). The equality subscheme is exactly the open and closed subset \(D(d)\) with this ring structure, including all nilpotents. Formula (UR-4) commutes with arbitrary further base change of \(B\). If \(\operatorname{Spec}(B)\) is nonempty and connected, then \(d\in\{0,1\}\); therefore equality at even one residue-field point forces \(d=1\) and \(f=g\) on the whole ring. Connectedness here does not require reducedness.

The finite diagonal condition has a real boundary. Take the already used tower \(K_n=\mathbf Q(2^{1/2^n})\) with compatible positive roots and \(L=\bigcup_nK_n\). The Eisenstein degrees \(2^n\) make this an infinite algebraic extension, and characteristic zero makes it separable. Every \(\alpha\in L\) satisfies a separable minimal polynomial \(F\); differentiating gives \(F'(\alpha)d\alpha=0\). The nonzero field element \(F'(\alpha)\) is a unit, hence \(\Omega_{L/\mathbf Q}=0\). But its diagonal admits no idempotent as above. If \(e=\sum a_i\otimes b_i\) existed, choose \(n\) with all \(b_i\in K_n\) and choose \(\alpha\in L\setminus K_n\). The equality \((\alpha\otimes1)e=(1\otimes\alpha)e\) would put the first side in \(L\otimes_{\mathbf Q}K_n\) and the second in \(L\otimes_{\mathbf Q}\alpha K_n\). These subspaces have zero intersection, because \(K_n\cap\alpha K_n=0\), and tensoring vector spaces over \(\mathbf Q\) preserves that direct sum. Both sides would be zero. Multiplication by \(1\otimes\alpha\) is invertible, forcing \(e=0\), contrary to \(\mu(e)=1\). Hence this formally unramified map has a diagonal ideal which is not finitely generated and no projective diagonal module. It is not a counterexample to the source's finite-type theorem.

\subsubsection{Report decisions and propagation}\label{algebra-proof-056-report-decisions-and-propagation}

Seven groups cover the twelve reports exactly once and give seven operations: paired article reports 12212/942; paired finite-presentation wording 12213/943; composition reference 944; local-to-global reference 945; paired maximal-ideal sentence 946/12214; paired plural verb 12215/947; and paired modal spelling 948/12216. The two reference repairs are mathematical corrections of the stated proof dependencies. The other changes preserve the exact hypotheses and conclusions as copyedits. No prior canonical correction lies in this source interval.

One zero-operation editorial group proves the exact determinant idempotent, the weaker finite-diagonal hypothesis and its converse via projectivity, arbitrary base change, the full equality factor for two maps, and the infinite separable-extension boundary. It receives the formal-unramified criterion from the preceding review and the exact field tower from group 1181, retaining every diagonal map and determinant coefficient. The stronger G-unramified choice already supplied by the finite-coefficient construction is recorded as existing proof content. All fuller arguments and consequences stay editorial; they do not replace the original exposition. Broader synthesis and its paper follow core completion.
